\section{Conclusiones}

La comparaci\'on entre ambas redes deja una idea simple: la centralidad no significa lo mismo en todos los grafos. En la minired del S\&P 500, las aristas distinguen funciones econ\'omicas y por eso la importancia se parte en dos preguntas distintas: qui\'en abastece y qui\'en compra. Ah\'i HITS funciona muy bien porque convierte esa diferencia en algo visible. La intermediaci\'on agrega otra capa al mostrar qu\'e empresas, como \textbf{AMZN} o \textbf{MSFT}, conectan rutas que de otro modo quedar\'ian m\'as sueltas.

En \textit{Star Wars Episode II}, en cambio, la red no separa roles direccionales. Lo que mide es qu\'e tan metido est\'a cada personaje en el centro del relato. Varias m\'etricas coinciden en el mismo trio: \textbf{Padm\'e}, \textbf{Obi-Wan} y \textbf{Anakin}. La diferencia m\'as interesante la aporta betweenness, porque ayuda a ver que \textbf{Obi-Wan} no solo participa mucho, sino que enlaza partes distintas de la historia.

En conjunto, el ejercicio muestra que un nodo puede ser central por razones muy diferentes: por vender a muchos, por comprar de nodos fuertes, por estar cerca del n\'ucleo espectral o por servir de puente. La mejor m\'etrica no es la m\'as famosa ni la m\'as sofisticada, sino la que mejor responde a la pregunta que uno quiere hacerle a la red.
